\chapter*{Abstract}
\addcontentsline{toc}{chapter}{Abstract}

Phase-field fracture methods replace a sharp crack by a regularized field and therefore require fine spatial resolution in the fracture process zone. This thesis investigates whether Abaqus-native error-guided remeshing can provide that resolution efficiently for a phase-field formulation implemented with user elements. The present draft deliberately concentrates on the Mode-I benchmark that must be understood before state transfer or more complex fracture configurations are assessed.

A fixed-mesh single-edge-crack reference was verified first. The project simulation with 15,192 finite elements reached a peak reaction force of $0.757778\,\mathrm{kN}$ at $0.005857\,\mathrm{mm}$. An ordinary-least-squares fit over the first 400 active increments gives the project-derived initial stiffness $K_0=137.945520\,\mathrm{kN/mm}$; this stiffness is not reported by Pandey and Kumar. The two-job adaptive workflow described by Pandey and Kumar was then reconstructed: a coarse \texttt{Job-1\_UEL} pre-analysis supplies the Abaqus stress discretization error indicator \texttt{MISESERI}, a \texttt{RemeshingRule} and \texttt{adaptiveRemesh} generate a new mesh, and the refined mesh is converted to the layered user-element representation for the fracture analysis.

A preprocessing defect was identified in the first 71,320-element fracture deck. A malformed \texttt{N\_BOTTOM} definition retained only 16 of 150 intended node identifiers, leaving 134 bottom nodes vertically unconstrained. Correct line wrapping restored the intended support. The corrected full-fracture calculation gave $K_0=137.820804\,\mathrm{kN/mm}$, $F_{\max}=0.745325\,\mathrm{kN}$, and $u(F_{\max})=0.005750\,\mathrm{mm}$. An independent elastic requalification gave $K_0=138.021013\,\mathrm{kN/mm}$. The stiffness anomaly was therefore an input-format defect rather than a consequence of adaptive-mesh physics.

For the project reconstruction using the published Listing-1 settings, a 2,906-element coarse mesh and a 1\% target tolerance produce 71,320 finite elements, whereas Pandey and Kumar report 13,941 for their proposed adaptive branch. Controlled sensitivity studies reproduce the expected trend that a larger target tolerance yields fewer elements, but they do not identify the authors' complete executed configuration. The exact count difference is consequently documented as a publication-information limitation, as accepted in the 17 September 2026 supervisor meeting. The next stage is to expose physically meaningful user-element energy quantities and establish global energy balance, complete force--displacement agreement, and spatial and temporal convergence before Mode-I state-transfer conservation is evaluated.
